\section{The source problem and construction parameters}
\label{sec:source}

Write \([k]=\{1,\ldots,k\}\). A \emph{Label Cover instance}
\(\Phi\) consists of two disjoint finite variable sets \(U,V\),
two nonempty finite alphabets \(\Sigma_U,\Sigma_V\), and an explicit
nonempty list of \(M\) constraint occurrences. Occurrence \(c\in[M]\)
has endpoints \(x_c\in U\), \(y_c\in V\), and a total map
\(\pi_c:\Sigma_U\to\Sigma_V\). A labeling \(A\) assigns each
variable a symbol from its side's alphabet. It satisfies occurrence
\(c\) when
\[
  \pi_c(A(x_c))=A(y_c).
\]
The value \(\val(\Phi)\) is the maximum fraction of satisfied
occurrences. List entries count with their multiplicity. Endpoints
and map tables are given explicitly in binary; unused variable
declarations can be removed. We require no regularity or expansion
condition on the constraint graph.

\begin{theorem}[Perfect-completeness Label Cover hardness]
\label{src:hardness}
For every fixed rational \(\sigma\in(0,1)\), there are nonempty
constant alphabets and a deterministic polynomial-time reduction
from \(\mathrm{3SAT}\) to instances of the form above such that
satisfiable formulas give \(\val(\Phi)=1\), and unsatisfiable
formulas give \(\val(\Phi)\leq\sigma\).
\end{theorem}

This is the standard consequence of the probabilistically checkable
proof (PCP) theorem~\cite{AroraSafra98,ALMSS98} and parallel
repetition~\cite{Raz1998}.
A quantitative formulation appears in the full version of Dinur
and Steurer~\cite[Theorem~6.2]{DinurSteurer2014}.
The precise unweighted, total-projection formulation is
recorded in \cite[Definitions~1.1--1.2 and Theorem~1.3]{DKKMS2025}.
There the constraints are even \(D\)-to-one for a suitable constant
\(D\), a property we do not use. The constants, including the
alphabets and \(D\), may depend on \(\sigma\). This theorem is the
only external hardness input to our proof.

\subsection{Parameters for the construction}

Fix the requested gap integer \(m\geq4\). We first construct and
analyze graphs for arbitrary fixed source alphabets and an even square
integer \(d\) satisfying
\begin{equation}
 d\geq4,\qquad \sqrt d>4m.
 \label{src:dimension-bounds}
\end{equation}
The number \(d\) will count the positions in the construction.
The soundness argument will impose further lower bounds on \(d\)
and determine the required source soundness. In
\Cref{sec:conclusion} we choose all these constants as functions of
\(m\), before invoking \Cref{src:hardness} to obtain the alphabets.

For the graph construction put
\begin{equation}
 a=\frac1{4m},\qquad p=2d+1,\qquad t=a\sqrt d.
 \label{src:construction-parameters}
\end{equation}
The graph will use a grid with \(p\) points and an adjacency threshold
\(2t\). These choices give
\[
 \frac dp<\frac12<\frac t2,\qquad p>2m.
\]
Both \(p\) and \(\sqrt d\) are integers, so \(t\) is rational.
For a source instance over fixed alphabets, set
\[
 q=\max\{|\Sigma_U|,|\Sigma_V|\},\qquad r=q^{d-1}.
\]
A query will involve at most \(d-1\) variables, and \(r\) bounds
its number of local labels. For fixed \(m,d\) and alphabets, every
parameter introduced here is constant with respect to the source
input size.
